\section{\textcolor{Salmon}{Zufriedenheit} \textcolor{red}{TODO: Titel}}
\label{ch:Diskussion:Zufriedenheit}

Übergang von Intuition über Gewohnheit\\
--> nur weil ein System intuitiv ist, müssen die Nutzer es nicht mögen\\

\noindent
irgendwas in Bezug auf \textit{Umstieg auf das neue System}\\
  --> gelingt es den Mitarbeiten\\
  --> was spricht dafür / dagegen\\
  --> neue Kollegen\\

% Allgemein
\noindent
Nutzer sind zufrieden damit\\
--> Es lässt sich zusammenfassend sagen, dass die Nutzer sehr zufrieden mit dem neuen Dashboard sind.\\
--> Prümper\\
--> Gibt es eine Diskrepanz zwischen Nutzen und Nutzbarkeit eines System (obj / subj)?\\
\textbf{Decken sich die Ergebnisse der Umfrage und des Tests?}\\
-	wie sich gemessene Werte und Umfrageergebnisse decken / bestätigen\\
-	Bewerten in jeder Nuance\\
  --> moderierende Variablen, Effizienz/Zeit, Effektivität/Genauigkeit, Umfrage\\
  --> alles in die Kategorien der ISO-Norm???\\

% Ausreißer
\noindent
außer: die 2 nicht signifikanten Fragen\\
  --> neu rechnen > doch signifikant\\

% Begründung
\noindent
warum? – die 2 mit der meisten Erfahrung und dem meisten Use in dem Bereich, haben allen anderen entgegengesprochen\\
  --> Gewohnheit\\

\noindent
Während auch DB\_use signifikant ist\\
  --> sind auch die Einzigen beiden, die Dashboard\_use mit 'nein' beantwortet haben\\

% weiteres
\noindent
- nur Vermutungen in Bezug auf das Alter möglich
